\scanpage{183}\mspage{38}
dans cette première approche pré-catégorie dérivée, par la connaissance du groupe fondamental $\pi_1$, \unsure{donc} des catégories des coeff. \unsure{constant|constants} tordus $=$ $\pi_1$-\unsure{modules|module}, les foncteurs $H^i(-,M)$ sur ceux-ci avec la structure de foncteur \deleted{\unsure{ab|ab.}, de \unsure{foncteur}} \inserted{cohomologique} relatif aux suites exactes, la \deleted{\unsure{struc}}structure de cup-produit etc -- reliés par certaines propriétés formelles. Une fois qu'on dispose du langage des cat. dérivées, la \inserted{sous catégorie de la catégorie dérivée des} complexes abéliens, \unsure{form-|formée|formés} \inserted{aux complexes de} faisceaux de coh. loc constante sur $X$, avec sa structure \unsure{triangulé|triangulée} et sa structure \deleted{\unsure{de|sou|m}}multiplicative $\overset{\mathbb{L}}{\otimes}$ (et éventuellement $\mathbb{R}\mathrm{Hom}$ \ldots) donnait un candidat plus satisfaisant pour espérer récupérer le type d'homotopie. Je n'ai guère d'idée si oui ou non ils permettent de le récupérer \inserted{en effet}, mais par contre n'ai guère de doute qu'en allant au bout de la ``linéarisation'', c'est à dire en partant dans le \uline{non abélien} et travaillant avec la $(n+1)$-catégorie (\unsure{sans|sa} \unsure{structures|structure} supplémentaires \unsure{dessus|disons}!) des $n$-champs loc. constants \inserted{\unsure{sur $X$}} pour tout $n$
% archive sheet number: 183
